% !PRINTED 133
\weberpage{133}

If the absolute values of $\varepsilon_1,\varepsilon_2,
\ldots$ lie below sufficiently small values, then we can prove first of the function $f_1(x)$
\begin{center}
\gesperrt{that it has no roots outside the regions $(\varrho),(\varrho'),(\varrho''),\ldots$,}
\end{center}
and secondly,
\begin{center}
\gesperrt{that the number of roots of $f_1(x)$ within $(\varrho)$ is exactly $a$, within $(\varrho')$ exactly $b$, within $(\varrho'')$ exactly $c$, and so on.}
\end{center}
In the last part of the theorem, however, it must be observed that, if $f_1(x)$ has multiple roots, these must be counted according to their multiplicity.

We construct, according to \S 31, IV, in the $x$-plane a circle with radius $R$ and with the origin as centre, which encloses the regions $(\varrho),(\varrho'),(\varrho''),\ldots$, so that outside this circle no more roots of $f_1(x)$ lie, and we denote by $G$ the region lying within this circle, but outside $(\varrho),(\varrho'),(\varrho''),\ldots$. It must also be remarked that $R$ can be assumed independently of $\varepsilon_1,\varepsilon_2,\ldots$, so long as a definite upper bound is fixed for the absolute values of these quantities.

% Exact unretouched crop from the registered Volume-I scan, printed Fig. 4.
\weberfigureimage{weber-b1-fig04.png}{0.42}

If now $x$ is a point of the region $G$, then the absolute value of $x-\alpha$ is greater than $\varrho$, and consequently, by (2),
\begin{equation}
 |f(x)|>\varrho^a\varrho'^b\varrho''^c\cdots .
\tag{5}
\end{equation}
If now $\alpha_1$ is a root of $f_1(x)$, then from (4) it follows that
\begin{equation}
 f(\alpha_1)=-\varphi(\alpha_1),
\tag{6}
\end{equation}
and $\alpha_1$ therefore cannot lie in the region $G$ if one takes care that everywhere within $G$
\begin{equation}
 |\varphi(x)|<\varrho^a\varrho'^b\varrho''^c\cdots,
\tag{7}
\end{equation}
which is always possible by sufficient diminution of the upper bound of $\varepsilon_1,\varepsilon_2,
\ldots$. Thus the first part of our assertion is proved, namely that, when the inequality (7) is satisfied, no roots of $f_1(x)$ lie outside the regions $(\varrho),(\varrho'),(\varrho''),\ldots$.
